% mathpreview built-in macro overrides.
%
% Best-effort stubs for macros that real papers reach for but which the
% preamble extractor cannot translate from the source — typically because
% the author's definition delegates to an @-internal helper (e.g.
% `\given` -> `\SV@given{\delimsize}` from svmacro.sty) that uses TeX
% primitives MathJax cannot expand. Without a stub, every equation that
% touches one of these commands fails to render.
%
% These are the lowest-priority entries in the cascade:
%   builtins (this file)
%     -> paper preamble macros
%       -> ~/.config/mathpreview/macros.tex
%         -> .mathpreview-macros.tex in the project root
%           -> --macros CLI flag
% Later definitions override earlier ones by name, so any author who
% provides a cleaner `\newcommand{\given}{...}` in their preamble — or
% in any of the user override files — wins automatically.

% \given / \st from svmacro.sty-style packages: spaced pipe used in
% `P(A \given B)` and `\{x \st x > 0\}`.
\newcommand{\given}{\,|\,}
\newcommand{\st}{\,|\,}

% \bm from the `bm` package. MathJax's `[tex]/boldsymbol` extension
% provides `\boldsymbol` but does not alias `\bm`, so author macros like
% `\newcommand{\E}{\bm{E}}` (svmacro.sty) break every equation they touch.
\newcommand{\bm}[1]{\boldsymbol{#1}}

% \underaccent from the `accents` package — closest pure-MathJax fit.
\newcommand{\underaccent}[2]{\underset{#1}{#2}}

% \xspace is a no-op in math contexts.
\newcommand{\xspace}{}

% \defeq (svmacro / mathtools-style) rendered as a stylized `:=`.
\newcommand{\defeq}{\stackrel{\scriptscriptstyle\mathrm{def}}{=}}

% Sidenote / annotation commands authors use for review comments.
% `\sidenote[opts]{text}` from a typical svmacro.sty / tcolorbox bridge.
% Stubbed to render as empty in math contexts — author-private review
% notes are noise in a live preview.
\newcommand{\sidenote}[2][]{}

% Edit-tracking commands from marktext-style packages.
\newcommand{\add}[1]{#1}
\newcommand{\remove}[1]{}
\newcommand{\highlight}[1]{#1}
\newcommand{\replace}[2]{#2}
